% ============================================================
%  Livrable Semaine 6 — Microfrontend React
%  Système de Gestion de Flotte de Véhicules
%  Master 1 GIL — 2025/2026
% ============================================================
\documentclass[12pt,a4paper]{report}

% ── Encodage & langue ──────────────────────────────────────
\usepackage[utf8]{inputenc}
\usepackage[T1]{fontenc}
\usepackage[french]{babel}

% ── Mise en page ──────────────────────────────────────────
\usepackage[top=2.5cm, bottom=2.5cm, left=2.5cm, right=2.5cm]{geometry}
\usepackage{setspace}
\setstretch{1.1}

% ── Typographie ───────────────────────────────────────────
\usepackage{lmodern}
\usepackage{microtype}

% ── Couleurs ──────────────────────────────────────────────
\usepackage[dvipsnames,svgnames,x11names,table]{xcolor}
\definecolor{titleblue}{HTML}{2C3E7A}
\definecolor{linkblue}{HTML}{2980B9}
\definecolor{lightblue}{HTML}{D6E8F7}
\definecolor{codebg}{HTML}{F4F6F8}
\definecolor{codeframe}{HTML}{BDC3C7}
\definecolor{keyword}{HTML}{2471A3}
\definecolor{comment}{HTML}{7F8C8D}
\definecolor{string}{HTML}{1E8449}

% ── Graphiques ────────────────────────────────────────────
\usepackage{graphicx}
\usepackage{tikz}
\usetikzlibrary{shapes.geometric, arrows.meta, positioning, backgrounds}

% ── Code ──────────────────────────────────────────────────
\usepackage{listings}
\usepackage{lstautogobble}

\lstdefinelanguage{TypeScript}{
  keywords={import,export,default,const,let,var,function,return,interface,type,
             class,extends,implements,new,if,else,for,while,async,await,from,of,
             in,true,false,null,undefined,void,string,number,boolean,React,
             useEffect,useState,useCallback,useRef,delete},
  sensitive=true,
  comment=[l]{//},
  morecomment=[s]{/*}{*/},
  morestring=[b]",
  morestring=[b]',
  morestring=[b]`,
}

\lstdefinestyle{code}{
  basicstyle=\ttfamily\small,
  keywordstyle=\color{keyword}\bfseries,
  commentstyle=\color{comment}\itshape,
  stringstyle=\color{string},
  numberstyle=\tiny\color{gray},
  numbers=left,
  stepnumber=1,
  numbersep=8pt,
  backgroundcolor=\color{codebg},
  frame=single,
  framerule=0.4pt,
  rulecolor=\color{codeframe},
  tabsize=2,
  showstringspaces=false,
  breaklines=true,
  breakatwhitespace=true,
  autogobble=true,
  xleftmargin=14pt,
  xrightmargin=4pt,
  aboveskip=8pt,
  belowskip=8pt,
}

\lstdefinestyle{shell}{
  basicstyle=\ttfamily\small,
  backgroundcolor=\color{codebg},
  frame=single,
  framerule=0.4pt,
  rulecolor=\color{codeframe},
  numbers=none,
  breaklines=true,
  autogobble=true,
  xleftmargin=14pt,
  aboveskip=8pt,
  belowskip=8pt,
}

\lstdefinestyle{sql}{
  language=SQL,
  basicstyle=\ttfamily\small,
  keywordstyle=\color{keyword}\bfseries,
  commentstyle=\color{comment}\itshape,
  stringstyle=\color{string},
  numbers=left,
  numberstyle=\tiny\color{gray},
  numbersep=8pt,
  backgroundcolor=\color{codebg},
  frame=single,
  framerule=0.4pt,
  rulecolor=\color{codeframe},
  tabsize=2,
  showstringspaces=false,
  breaklines=true,
  autogobble=true,
  xleftmargin=14pt,
  aboveskip=8pt,
  belowskip=8pt,
}

% ── Tables ────────────────────────────────────────────────
\usepackage{booktabs}
\usepackage{tabularx}
\usepackage{longtable}
\usepackage{multirow}
\usepackage{array}
\newcolumntype{L}[1]{>{\raggedright\arraybackslash}p{#1}}
\newcolumntype{C}[1]{>{\centering\arraybackslash}p{#1}}

% ── Listes ────────────────────────────────────────────────
\usepackage{enumitem}
\setlist[itemize]{noitemsep, topsep=3pt, leftmargin=*}
\setlist[enumerate]{noitemsep, topsep=3pt, leftmargin=*}

% ── Liens ─────────────────────────────────────────────────
\usepackage[colorlinks=true, linkcolor=titleblue, urlcolor=linkblue, citecolor=linkblue]{hyperref}

% ── En-têtes / pieds ──────────────────────────────────────
\usepackage{fancyhdr}
\pagestyle{fancy}
\fancyhf{}
\fancyhead[L]{\small Gestion de Flotte de Véhicules}
\fancyhead[R]{\small Master 1 GIL — 2025/2026}
\fancyfoot[C]{\small\thepage}
\renewcommand{\headrulewidth}{0.4pt}

% ── Titres de chapitres ───────────────────────────────────
\usepackage{titlesec}
\titleformat{\chapter}[display]
  {\normalfont\huge\bfseries\color{titleblue}}
  {\chaptertitlename\ \thechapter}{12pt}{\Huge}
\titlespacing*{\chapter}{0pt}{-20pt}{20pt}

\titleformat{\section}{\large\bfseries\color{titleblue}}{}{0em}{\thesection\quad}
\titleformat{\subsection}{\normalsize\bfseries\color{titleblue}}{}{0em}{\thesubsection\quad}
\titleformat{\subsubsection}{\normalsize\bfseries}{}{0em}{\thesubsubsection\quad}

% ── Numérotation ──────────────────────────────────────────
\setcounter{secnumdepth}{3}
\setcounter{tocdepth}{3}

% ── Utilitaires ───────────────────────────────────────────
\usepackage{parskip}
\usepackage{caption}
\captionsetup{font=small, labelfont=bf}

% ============================================================
\begin{document}
\thispagestyle{empty}

% ── Page de garde ─────────────────────────────────────────
\begin{titlepage}
  \centering
  \vspace*{2cm}

  \rule{\linewidth}{1.5pt}\\[0.4cm]
  {\Huge\bfseries Rapport Technique}\\[0.3cm]
  {\Large\color{titleblue} Microfrontend React — Interface Utilisateur}\\[0.2cm]
  {\normalsize Intégration SSO Keycloak, RBAC, Leaflet GPS et Tests E2E Cypress}\\[0.1cm]
  \rule{\linewidth}{1.5pt}\\[1.5cm]

  {\Large Système de Gestion de Flotte de Véhicules}\\[0.5cm]
  {\large Master 1 GIL — 2025/2026}\\[2cm]

  \begin{minipage}{0.6\textwidth}
    \begin{tabular}{|l l|}
      \hline
      \multicolumn{2}{|c|}{\textbf{Informations du projet}} \\
      \hline
      \textbf{Auteurs :}            & Kherbachi Rima \& Mouhoubi Leiticia \\
      \textbf{Formation :}          & Master 1 GIL \\
      \textbf{Université :}         & Université de Rouen Normandie \\
      \textbf{Encadrants :}         & Lydia et Luc \\
      \textbf{Année :}              & 2025 – 2026 \\
      \textbf{Semaine documentée :} & 6 \\
      \hline
    \end{tabular}
  \end{minipage}

  \vfill
  {\large Mars – Avril 2026}
\end{titlepage}

% ── Table des matières ────────────────────────────────────
\newpage
\tableofcontents
\newpage

% ============================================================
\chapter{Introduction et Architecture}
% ============================================================

\section{Présentation du livrable}

La semaine 6 est consacrée au développement du \textbf{microfrontend React}, l'interface
graphique unifiée du système de gestion de flotte. Elle consolide l'ensemble des microservices
développés lors des semaines précédentes en les exposant à travers une \textbf{Single Page Application}
(SPA) sécurisée, accessible selon les rôles métier.

\begin{center}
\begin{tabular}{|l|l|l|}
  \hline
  \rowcolor{lightblue}
  \textbf{Technologie} & \textbf{Version} & \textbf{Rôle} \\
  \hline
  React           & 19.2   & Framework UI (composants fonctionnels + hooks) \\
  TypeScript      & $\sim$6.0 & Typage statique \\
  Vite            & 8.0    & Bundler / dev server (HMR) \\
  keycloak-js     & 26.2   & Client SSO OIDC / PKCE S256 \\
  react-leaflet   & 5.0    & Carte interactive GPS (React 19 compatible) \\
  Leaflet         & 1.9    & Moteur cartographique OpenStreetMap \\
  Recharts        & 3.8    & Graphiques SVG (vitesse, activité) \\
  socket.io-client & 4.8   & WebSocket temps réel (positions GPS) \\
  Axios           & 1.15   & Client HTTP avec intercepteur JWT \\
  lucide-react    & 1.8    & Icônes SVG \\
  Cypress         & 15.13  & Tests end-to-end \\
  \hline
\end{tabular}
\captionof{table}{Pile technologique du microfrontend}
\end{center}

\section{Structure du projet}

\begin{lstlisting}[style=shell, caption=Arborescence de \texttt{frontend-legacy/src/}]
src/
├── auth.ts                  # Initialisation Keycloak (PKCE S256)
├── api.ts                   # Axios + intercepteur JWT auto-refresh
├── App.tsx                  # Router + Sidebar + RBAC routing
├── App.css                  # Theme sombre + variables CSS
├── index.css                # Reset global + scrollbar personnalisee
├── main.tsx                 # Entry point (initKeycloak + CSS Leaflet)
├── components/
│   ├── ProtectedRoute.tsx   # Garde RBAC par role Keycloak
│   ├── StatusBadge.tsx      # Badge colore par statut metier
│   └── Modal.tsx            # Dialog generique (ESC + overlay)
└── pages/
    ├── Dashboard.tsx        # KPI Socket.IO + Recharts
    ├── Vehicules.tsx        # CRUD vehicules + filtres
    ├── Conducteurs.tsx      # CRUD conducteurs (cards)
    ├── Localisation.tsx     # Carte Leaflet temps reel
    ├── Maintenance.tsx      # Interventions + workflow statuts
    └── Unauthorized.tsx     # Page 403 RBAC
\end{lstlisting}

\section{Diagramme de flux}

\begin{center}
\begin{tikzpicture}[
  node distance=1.4cm,
  box/.style={rectangle, draw=titleblue!70, fill=lightblue!60, rounded corners=4pt,
              minimum width=3cm, minimum height=0.8cm, font=\small\bfseries, align=center},
  svc/.style={rectangle, draw=gray!60, fill=gray!10, rounded corners=3pt,
              minimum width=2.6cm, minimum height=0.7cm, font=\small, align=center},
  arr/.style={->, >=Stealth, thick, color=titleblue},
  darr/.style={->, >=Stealth, dashed, color=gray!70},
]

\node[box] (browser) {Navigateur};
\node[box, right=2.5cm of browser] (kc) {Keycloak\\\small :9080};
\node[box, below=1.4cm of browser] (app) {React App (Vite)};
\node[svc, below left=1.2cm and 0.5cm of app] (pages) {Pages (5)};
\node[box, right=2.8cm of app] (gw) {API Gateway\\\small :4000};
\node[box, below=1.2cm of gw] (ws) {Socket.IO\\\small :3002};

\draw[arr] (browser) -- node[above,font=\scriptsize]{OIDC/PKCE} (kc);
\draw[arr] (browser) -- (app);
\draw[arr] (app) -- node[above,font=\scriptsize]{JWT Bearer} (gw);
\draw[arr] (app) -- node[right,font=\scriptsize]{WebSocket} (ws);
\draw[darr] (app) -- (pages);
\end{tikzpicture}
\end{center}

% ============================================================
\chapter{Authentification SSO — Keycloak}
% ============================================================

\section{Configuration Keycloak}

L'authentification repose sur \textbf{Keycloak 26} avec le flux
Authorization Code + \textbf{PKCE S256}, garantissant qu'aucun secret client n'est
exposé dans le code frontend.

\begin{center}
\begin{tabular}{|l|l|}
  \hline
  \rowcolor{lightblue}
  \textbf{Paramètre} & \textbf{Valeur} \\
  \hline
  Realm     & \texttt{fleet-management} \\
  Client ID & \texttt{fleet-api-gateway} \\
  URL       & \texttt{http://localhost:9080} \\
  Flow      & Authorization Code + PKCE S256 \\
  onLoad    & \texttt{login-required} \\
  \hline
\end{tabular}
\captionof{table}{Paramètres de configuration Keycloak}
\end{center}

\begin{lstlisting}[style=code, language=TypeScript, caption=\texttt{src/auth.ts} — Initialisation Keycloak]
import Keycloak from 'keycloak-js';

const keycloak = new Keycloak({
  url:      'http://localhost:9080',
  realm:    'fleet-management',
  clientId: 'fleet-api-gateway',
});

export const initKeycloak = (onAuthenticatedCallback: () => void) => {
  keycloak
    .init({ onLoad: 'login-required',
            checkLoginIframe: false,
            pkceMethod: 'S256' })
    .then(authenticated => {
      if (!authenticated) console.warn('Non authentifie');
      onAuthenticatedCallback();
    });
};

export default keycloak;
\end{lstlisting}

\section{Intercepteur JWT Axios}

Chaque requête HTTP est automatiquement enrichie avec le token Bearer.
En cas de réponse 401, le token est renouvelé silencieusement via
\texttt{keycloak.updateToken()}, sans interruption pour l'utilisateur.

\begin{lstlisting}[style=code, language=TypeScript, caption=\texttt{src/api.ts} — Intercepteur JWT avec auto-refresh]
const api = axios.create({ baseURL: 'http://localhost:4000' });

api.interceptors.request.use(async (config) => {
  await keycloak.updateToken(70); // Renouvelle si expiration < 70s
  config.headers.Authorization = `Bearer ${keycloak.token}`;
  return config;
});

api.interceptors.response.use(
  res => res,
  async (error) => {
    if (error.response?.status === 401) {
      await keycloak.updateToken(-1); // Force le renouvellement
    }
    return Promise.reject(error);
  }
);
\end{lstlisting}

% ============================================================
\chapter{Contrôle d'Accès Basé sur les Rôles (RBAC)}
% ============================================================

\section{Rôles et permissions}

Quatre rôles Keycloak de niveau \textit{realm} sont définis. Chaque rôle
détermine les pages accessibles et les actions autorisées.

\begin{center}
\begin{tabular}{|l|l|l|}
  \hline
  \rowcolor{lightblue}
  \textbf{Rôle} & \textbf{Permissions} & \textbf{Pages accessibles} \\
  \hline
  \texttt{admin}
    & Accès complet, tous CRUD
    & Toutes \\
  \hline
  \texttt{manager}
    & Gestion véhicules + conducteurs
    & Dashboard, Véhicules, Conducteurs, \newline Localisation, Maintenance \\
  \hline
  \texttt{technicien}
    & Lecture véhicules, gestion maintenance
    & Dashboard, Véhicules, Localisation, \newline Maintenance \\
  \hline
  \texttt{utilisateur}
    & Lecture seule
    & Dashboard, Véhicules, Localisation \\
  \hline
\end{tabular}
\captionof{table}{Matrice des rôles et permissions}
\end{center}

\section{Composant ProtectedRoute}

Le composant \texttt{ProtectedRoute} encapsule la logique de garde. Il interroge
\texttt{keycloak.hasRealmRole()} pour chaque rôle autorisé, et redirige vers
\texttt{/unauthorized} (page 403) si aucun ne correspond.

\begin{lstlisting}[style=code, language=TypeScript, caption=\texttt{src/components/ProtectedRoute.tsx}]
import { Navigate } from 'react-router-dom';
import keycloak from '../auth';

interface ProtectedRouteProps {
  children: React.ReactNode;
  roles: string[];
}

const ProtectedRoute = ({ children, roles }: ProtectedRouteProps) => {
  const hasAccess = roles.some(role => keycloak.hasRealmRole(role));
  if (!hasAccess) return <Navigate to="/unauthorized" replace />;
  return <>{children}</>;
};

export default ProtectedRoute;
\end{lstlisting}

\section{Routage RBAC dans App.tsx}

\begin{lstlisting}[style=code, language=TypeScript, caption=Extrait des routes protegees dans \texttt{App.tsx}]
<Routes>
  <Route path="/"             element={<Dashboard />} />
  <Route path="/unauthorized" element={<Unauthorized />} />

  <Route path="/vehicules" element={
    <ProtectedRoute roles={['admin','manager','technicien','utilisateur']}>
      <Vehicules />
    </ProtectedRoute>
  } />
  <Route path="/conducteurs" element={
    <ProtectedRoute roles={['admin','manager']}>
      <Conducteurs />
    </ProtectedRoute>
  } />
  <Route path="/localisation" element={
    <ProtectedRoute roles={['admin','manager','technicien','utilisateur']}>
      <Localisation />
    </ProtectedRoute>
  } />
  <Route path="/maintenance" element={
    <ProtectedRoute roles={['admin','manager','technicien']}>
      <Maintenance />
    </ProtectedRoute>
  } />
</Routes>
\end{lstlisting}

% ============================================================
\chapter{Pages et Fonctionnalités}
% ============================================================

\section{Tableau de bord (Dashboard)}

Le Dashboard agrège les métriques en temps réel reçues via \textbf{Socket.IO}
depuis le \texttt{localisation-service} (port 3002).

\begin{center}
\begin{tabular}{|l|l|}
  \hline
  \rowcolor{lightblue}
  \textbf{Indicateur} & \textbf{Source} \\
  \hline
  Nombre de véhicules actifs   & Socket.IO — positions des 5 dernières minutes \\
  Vitesse moyenne (km/h)       & Calcul sur les 20 derniers événements \\
  Total événements             & Compteur cumulé des positions reçues \\
  Graphique de vitesse         & Recharts \texttt{AreaChart} SVG \\
  Flux d'activité              & Liste des 10 dernières positions horodatées \\
  \hline
\end{tabular}
\captionof{table}{Indicateurs du tableau de bord}
\end{center}

\section{Gestion des Véhicules}

Interface CRUD complète exposant le \texttt{vehicule-service} via l'API Gateway.

\begin{itemize}
  \item \textbf{Tableau} avec colonnes : Immatriculation, Marque/Modèle, Année, Statut, Actif, Actions
  \item \textbf{Filtres} par statut : Tous / DISPONIBLE / EN\_COURSE / EN\_MAINTENANCE
  \item \textbf{Changement de statut} inline via \texttt{<select>} (admin/manager uniquement)
  \item \textbf{Modal création} : immatriculation, marque, modèle, année
  \item \textbf{Modal édition} : modification de toutes les propriétés
  \item \textbf{Suppression logique} : \texttt{DELETE /vehicules/\{id\}} — désactivation du champ \texttt{actif}
\end{itemize}

\begin{center}
\begin{tabular}{|l|l|l|}
  \hline
  \rowcolor{lightblue}
  \textbf{Méthode} & \textbf{Endpoint} & \textbf{Action} \\
  \hline
  GET    & \texttt{/vehicules}                  & Liste des véhicules \\
  POST   & \texttt{/vehicules}                  & Créer un véhicule \\
  PUT    & \texttt{/vehicules/\{id\}}            & Mettre à jour \\
  PATCH  & \texttt{/vehicules/\{id\}/statut}     & Changer le statut \\
  DELETE & \texttt{/vehicules/\{id\}}            & Désactivation logique \\
  \hline
\end{tabular}
\captionof{table}{Endpoints REST — vehicule-service}
\end{center}

\section{Gestion des Conducteurs}

Interface en \textbf{grille de cartes} avec avatar généré à partir des initiales.
Chaque carte affiche les informations du permis avec un code couleur sur la date d'expiration :

\begin{itemize}
  \item \textcolor{green!60!black}{Vert} si l'expiration est à plus de 30 jours
  \item \textcolor{orange}{Orange} si l'expiration est dans moins de 30 jours
  \item \textcolor{red}{Rouge} si le permis est expiré
\end{itemize}

Une barre de recherche filtre les conducteurs en temps réel (côté client) par nom ou email.

\section{Carte de Localisation (Leaflet)}

\begin{lstlisting}[style=code, language=TypeScript, caption=Fix icone Leaflet avec Vite (bundler ES modules)]
// Les bundlers ES cassent le mecanisme interne de Leaflet
delete (L.Icon.Default.prototype as unknown as Record<string, unknown>)
  ._getIconUrl;
L.Icon.Default.mergeOptions({
  iconUrl: 'https://unpkg.com/leaflet@1.9.4/dist/images/marker-icon.png',
  iconRetinaUrl:
    'https://unpkg.com/leaflet@1.9.4/dist/images/marker-icon-2x.png',
  shadowUrl:
    'https://unpkg.com/leaflet@1.9.4/dist/images/marker-shadow.png',
});
\end{lstlisting}

\begin{center}
\begin{tabular}{|l|l|}
  \hline
  \rowcolor{lightblue}
  \textbf{Fonctionnalité} & \textbf{Implémentation} \\
  \hline
  Fond de carte          & OpenStreetMap (tuiles libres, sans clé API) \\
  Centre initial         & Rouen (lat 49.4431, lng 1.0993) \\
  Marqueurs colorés      & \texttt{DivIcon} HTML — vert/bleu/orange par statut \\
  Popup au clic          & Immatriculation, statut, vitesse, heure \\
  Zones géofencing       & Composant \texttt{<Circle>} avec remplissage semi-transparent \\
  Ajustement automatique & Hook \texttt{useMap()} + \texttt{fitBounds} sur positions \\
  Temps réel             & Socket.IO — mise à jour des marqueurs à chaque position \\
  Indicateur connexion   & Badge connecté/hors ligne (état WebSocket) \\
  \hline
\end{tabular}
\captionof{table}{Fonctionnalités de la carte Leaflet}
\end{center}

\begin{lstlisting}[style=code, language=TypeScript, caption=Composant FitBounds — ajustement automatique du zoom]
function FitBounds({ positions }: { positions: VehiculePosition[] }) {
  const map = useMap();
  useEffect(() => {
    if (positions.length === 0) return;
    const bounds = L.latLngBounds(
      positions.map(p => [p.latitude, p.longitude])
    );
    map.fitBounds(bounds, { padding: [40, 40] });
  }, []);
  return null;
}
\end{lstlisting}

\section{Gestion Maintenance}

\begin{center}
\begin{tabular}{|l|l|l|}
  \hline
  \rowcolor{lightblue}
  \textbf{Méthode} & \textbf{Endpoint} & \textbf{Action} \\
  \hline
  GET   & \texttt{/maintenance/interventions}              & Liste des interventions \\
  POST  & \texttt{/maintenance/interventions}              & Créer une intervention \\
  PATCH & \texttt{/maintenance/interventions/\{id\}/statut} & Changer le statut \\
  \hline
\end{tabular}
\captionof{table}{Endpoints REST — maintenance-service}
\end{center}

Le workflow des statuts est piloté par des boutons contextuels :

\begin{center}
\begin{tikzpicture}[
  node distance=2cm,
  s/.style={rectangle, draw=titleblue!60, fill=lightblue!40, rounded corners=3pt,
            minimum width=2.4cm, minimum height=0.7cm, font=\small\bfseries},
  arr/.style={->, >=Stealth, thick, color=titleblue},
]
\node[s] (p) {PLANIFIEE};
\node[s, right=2.2cm of p] (ec) {EN\_COURS};
\node[s, right=2.2cm of ec] (t) {TERMINEE};
\node[s, below=1cm of ec] (a) {ANNULEE};

\draw[arr] (p)  -- node[above,font=\scriptsize]{Démarrer} (ec);
\draw[arr] (ec) -- node[above,font=\scriptsize]{Terminer} (t);
\draw[arr] (p)  -- node[left,font=\scriptsize]{Annuler} (a);
\draw[arr] (ec) -- node[right,font=\scriptsize]{Annuler} (a);
\end{tikzpicture}
\end{center}

% ============================================================
\chapter{Composants Partagés}
% ============================================================

\section{StatusBadge}

Composant générique qui affiche un badge coloré pour tout statut métier.
Il accepte une chaîne ou un booléen et dispose d'un fallback gris pour les
valeurs inconnues.

\begin{center}
\begin{tabular}{|l|l|l|}
  \hline
  \rowcolor{lightblue}
  \textbf{Clé} & \textbf{Label affiché} & \textbf{Couleur} \\
  \hline
  \texttt{DISPONIBLE}      & Disponible  & Vert   \texttt{\#10b981} \\
  \texttt{EN\_COURSE}      & En Course   & Bleu   \texttt{\#3b82f6} \\
  \texttt{EN\_MAINTENANCE} & Maintenance & Orange \texttt{\#f59e0b} \\
  \texttt{PLANIFIEE}       & Planifiée   & Violet \texttt{\#6366f1} \\
  \texttt{EN\_COURS}       & En Cours    & Bleu   \texttt{\#3b82f6} \\
  \texttt{TERMINEE}        & Terminée    & Vert   \texttt{\#10b981} \\
  \texttt{ANNULEE}         & Annulée     & Rouge  \texttt{\#ef4444} \\
  \texttt{PENDING}         & En attente  & Orange \texttt{\#f59e0b} \\
  \texttt{VALIDE}          & Validée     & Vert   \texttt{\#10b981} \\
  \texttt{REJETEE}         & Rejetée     & Rouge  \texttt{\#ef4444} \\
  \texttt{true}            & Actif       & Vert   \texttt{\#10b981} \\
  \texttt{false}           & Inactif     & Rouge  \texttt{\#ef4444} \\
  \hline
\end{tabular}
\captionof{table}{Configuration du composant StatusBadge}
\end{center}

\section{Modal}

Dialog générique accessible :

\begin{itemize}
  \item Fermeture par touche \textbf{Échap} (listener sur \texttt{document})
  \item Fermeture par \textbf{clic sur l'overlay} semi-transparent
  \item Fermeture par bouton \textbf{✕} dans l'en-tête
  \item Largeur configurable via la prop \texttt{width} (défaut : 520px)
  \item Attributs \texttt{data-testid} sur tous les éléments pour Cypress
  \item Fond avec \texttt{backdrop-filter: blur(4px)}
\end{itemize}

\section{Design System — Thème Sombre}

\begin{lstlisting}[style=shell, caption=Variables CSS du design system (\texttt{App.css})]
:root {
  --primary:      #6366f1;              /* Indigo     */
  --bg-dark:      #0f172a;              /* Fond app   */
  --bg-sidebar:   #1e293b;              /* Fond sidebar */
  --glass:        rgba(255,255,255,0.05);
  --border-glass: rgba(255,255,255,0.1);
  --text-main:    #f8fafc;
  --text-muted:   #94a3b8;
}
\end{lstlisting}

Le layout principal utilise \texttt{display: grid} avec deux colonnes :
\texttt{240px} pour la sidebar et \texttt{1fr} pour le contenu.

% ============================================================
\chapter{Tests End-to-End avec Cypress}
% ============================================================

\section{Stratégie de tests}

\begin{center}
\begin{tabular}{|l|C{1.5cm}|l|}
  \hline
  \rowcolor{lightblue}
  \textbf{Fichier} & \textbf{Tests} & \textbf{Scénarios couverts} \\
  \hline
  \texttt{01\_auth.cy.ts}          & 9  & Sidebar, topbar, navigation, liens \\
  \texttt{02\_dashboard.cy.ts}     & 4  & KPI, Recharts SVG, WebSocket \\
  \texttt{03\_vehicules.cy.ts}     & 10 & CRUD, filtres, modal (ouverture/fermeture/soumission) \\
  \texttt{04\_conducteurs.cy.ts}   & 9  & Grille, recherche, CRUD, cartes \\
  \texttt{05\_localisation.cy.ts}  & 7  & Carte Leaflet, WebSocket, sélection véhicule \\
  \texttt{06\_maintenance.cy.ts}   & 13 & KPI, filtres, workflow statuts, CRUD \\
  \texttt{07\_rbac.cy.ts}          & 7  & Page 403, navigation active, routes inconnues \\
  \texttt{08\_status\_badge.cy.ts} & 8  & Labels, couleurs, badges par statut \\
  \hline
  \textbf{Total}                   & \textbf{67} & \\
  \hline
\end{tabular}
\captionof{table}{Couverture de la suite Cypress — 67 tests end-to-end}
\end{center}

\section{Commandes personnalisées}

Deux commandes sont définies dans \texttt{cypress/support/commands.ts} :

\begin{lstlisting}[style=code, language=TypeScript, caption=Commandes personnalisees Cypress]
// cy.login(role) — stub Keycloak sans serveur SSO
Cypress.Commands.add('login',
  (role: 'admin' | 'manager' | 'technicien' | 'utilisateur' = 'admin') => {
    cy.window().then((win) => {
      (win as any).__KC_STUB__ = {
        authenticated: true,
        token: 'fake-jwt-token',
        tokenParsed: { preferred_username: `test-${role}` },
        hasRealmRole: (r: string) => r === role,
        logout: () => {},
        updateToken: () => Promise.resolve(true),
      };
    });
  }
);

// cy.mockApi() — intercept REST avec fixtures JSON
Cypress.Commands.add('mockApi', () => {
  cy.intercept('GET', '**/vehicules',
    { fixture: 'vehicles.json' }).as('getVehicules');
  cy.intercept('GET', '**/conducteurs',
    { fixture: 'conducteurs.json' }).as('getConducteurs');
  cy.intercept('GET', '**/maintenance/interventions',
    { fixture: 'maintenance.json' }).as('getInterventions');
});
\end{lstlisting}

\section{Exemples de cas de tests}

\begin{lstlisting}[style=code, language=TypeScript, caption=Tests CRUD Vehicules (extrait)]
it('filtre par statut DISPONIBLE', () => {
  cy.contains('button', 'DISPONIBLE').click();
  cy.contains('AB-123-CD').should('be.visible');
  cy.contains('IJ-789-KL').should('not.exist');
});

it('ouvre le modal de creation via le bouton Nouveau', () => {
  cy.get('[data-testid="btn-create-vehicule"]').click();
  cy.get('[data-testid="modal-content"]').should('be.visible');
  cy.contains('Nouveau vehicule').should('be.visible');
});

it('ferme le modal avec la touche Echap', () => {
  cy.get('[data-testid="btn-create-vehicule"]').click();
  cy.get('body').type('{esc}');
  cy.get('[data-testid="modal-content"]').should('not.exist');
});

it('soumet le formulaire de creation', () => {
  cy.intercept('POST', '**/vehicules',
    { statusCode: 201, body: { id: 'veh-new' } }).as('createVehicule');

  cy.get('[data-testid="btn-create-vehicule"]').click();
  cy.get('[data-testid="input-immatriculation"]').type('QR-999-ST');
  cy.get('[data-testid="input-marque"]').type('Toyota');
  cy.get('[data-testid="input-modele"]').type('Yaris');
  cy.get('[data-testid="input-annee"]').type('2024');
  cy.get('[data-testid="create-vehicule-form"]').submit();

  cy.wait('@createVehicule').its('request.body')
    .should('deep.include', { immatriculation: 'QR-999-ST' });
});
\end{lstlisting}

\section{Commandes d'exécution}

\begin{lstlisting}[style=shell, caption=Lancement des tests Cypress]
# Mode interactif (interface graphique)
npx cypress open

# Mode headless (CI/CD)
npx cypress run

# Fichier specifique
npx cypress run --spec "cypress/e2e/03_vehicules.cy.ts"

# Avec rapport JUnit (integration CI)
npx cypress run --reporter junit
\end{lstlisting}

% ============================================================
\chapter{Déploiement et Points Techniques}
% ============================================================

\section{Variables d'environnement}

\begin{center}
\begin{tabular}{|l|l|C{1.5cm}|}
  \hline
  \rowcolor{lightblue}
  \textbf{Variable} & \textbf{Valeur par défaut} & \textbf{Requis} \\
  \hline
  \texttt{VITE\_API\_URL}        & \texttt{http://localhost:4000} & Non \\
  \texttt{VITE\_WS\_URL}         & \texttt{http://localhost:3002} & Non \\
  \texttt{VITE\_KC\_URL}         & \texttt{http://localhost:9080} & Oui \\
  \texttt{VITE\_KC\_REALM}       & \texttt{fleet-management}      & Oui \\
  \texttt{VITE\_KC\_CLIENT}      & \texttt{fleet-api-gateway}     & Oui \\
  \hline
\end{tabular}
\captionof{table}{Variables d'environnement Vite du frontend}
\end{center}

\section{Difficultés rencontrées et solutions}

\begin{center}
\begin{tabular}{|L{5cm}|L{8.5cm}|}
  \hline
  \rowcolor{lightblue}
  \textbf{Difficulté} & \textbf{Solution apportée} \\
  \hline
  \texttt{react-leaflet@4.x} incompatible React 19
    & Migration vers \texttt{react-leaflet@5.0} (React 19 natif) \\
  \hline
  \texttt{index.css} Vite fixait \texttt{\#root} à 1126px
    & Remplacement par reset minimal : \texttt{width: 100\%; height: 100vh} \\
  \hline
  Carte Leaflet non stylisée (CSS absent)
    & \texttt{import 'leaflet/dist/leaflet.css'} dans \texttt{main.tsx} avant l'app CSS \\
  \hline
  Icônes Leaflet manquantes avec Vite (bundler ES)
    & Suppression de \texttt{\_getIconUrl} + \texttt{mergeOptions} avec URLs CDN \\
  \hline
  Cypress ne peut pas contourner Keycloak
    & Commande \texttt{cy.login(role)} injectant un stub sur \texttt{window} \\
  \hline
  TypeScript 6 strict : cast Leaflet
    & \texttt{as unknown as Record<string, unknown>} pour le prototype \\
  \hline
\end{tabular}
\captionof{table}{Difficultés rencontrées et solutions apportées}
\end{center}

\section{Construction et démarrage}

\begin{lstlisting}[style=shell, caption=Scripts npm disponibles]
# Installation des dependances
npm install

# Serveur de developpement (http://localhost:5173)
npm run dev

# Build de production (TypeScript + Vite)
npm run build

# Previsualisation du build
npm run preview

# Lint ESLint
npm run lint
\end{lstlisting}

% ============================================================
\chapter*{Conclusion}
\addcontentsline{toc}{chapter}{Conclusion}
% ============================================================

La semaine 6 marque la livraison du \textbf{microfrontend React}, interface unifiée
du système FleetX. Elle consolide cinq semaines de développement microservices
en les exposant à travers une SPA professionnelle, sécurisée et testée.

\textbf{Réalisations principales :}
\begin{itemize}
  \item SPA React 19 complète avec thème sombre et design system cohérent.
  \item SSO Keycloak 26 (PKCE S256, \textit{login-required}) — zéro secret côté client.
  \item RBAC : 4 rôles, routes protégées, page 403.
  \item Dashboard temps réel Socket.IO + graphiques Recharts.
  \item Carte Leaflet avec marqueurs GPS colorés, géofencing et ajustement automatique.
  \item CRUD Véhicules (tableau + filtres + modals + workflow statuts).
  \item CRUD Conducteurs (grille de cartes + recherche + indicateur permis).
  \item Gestion Maintenance (workflow PLANIFIEE → EN\_COURS → TERMINEE/ANNULEE).
  \item \textbf{67 tests Cypress end-to-end} répartis en 8 fichiers.
  \item Composants partagés réutilisables : StatusBadge, Modal, ProtectedRoute.
\end{itemize}

\vspace{0.4cm}
\textbf{Avancement global du projet : 6/7 semaines — 86\,\%}

\end{document}
